
\documentclass[10pt,a4paper]{article}

% -----------------------------------------------------
% PACKAGES
% -----------------------------------------------------
\usepackage[margin=0.82in]{geometry}
\usepackage{hyperref}
\usepackage{enumitem}
\usepackage{graphicx}
\usepackage{xcolor}
\usepackage{titling}
\usepackage{titlesec}
\usepackage{times}

% -----------------------------------------------------
% DOCUMENT SETTINGS
% -----------------------------------------------------

\hypersetup{
    colorlinks=true,
    linkcolor=black,
    urlcolor=blue
}

\setlength{\parindent}{0pt}
\setlength{\parskip}{4pt}

\setlist[itemize]{
    leftmargin=0.55cm,
    itemsep=1pt,
    topsep=2pt,
    parsep=0pt
}

\setlist[enumerate]{
    leftmargin=0.55cm,
    itemsep=1pt,
    topsep=2pt,
    parsep=0pt
}

% Compact section spacing
\titlespacing*{\section}{0pt}{9pt}{3pt}
\titlespacing*{\subsection}{0pt}{6pt}{2pt}

% Section hints
\newcommand{\BvrHint}[1]{
    {\normalfont\normalsize\itshape\hfill #1}
}

% -----------------------------------------------------
% TITLE
% -----------------------------------------------------

\title{
    \vspace{-1.0cm}
    \textbf{STELLA: Solar Temporal Event Learning \& Likelihood Assessment}
}

\posttitle{
    \par\end{center}
    \begin{center}
    \normalsize Project Proposal for UCS503P\\
    Thapar Institute of Engineering and Technology
    \end{center}
    \vskip 0.25em
}

\author{
\begin{tabular}{cc}
    \textbf{Aditya Kumar} & \textbf{Pankhil Gupta} \\
    1024170314 & 1024170322 \\[12pt]
    \textbf{Hitanshi Antil} & \textbf{Ashish Kumar} \\
    1024170327 & 1024170341
\end{tabular}
}

\date{\today}

% -----------------------------------------------------
% DOCUMENT
% -----------------------------------------------------

\begin{document}

\maketitle

% -----------------------------------------------------
% 1. HIGHER-ORDER GOAL
% -----------------------------------------------------

\section{Higher-order goal%
  \BvrHint{\ldots secondary framing}}

This project aims to improve space-weather preparedness by reducing the time between
the emergence of a significant solar event and an actionable warning. While solar
flare forecasting is a broader scientific problem, the immediate engineering
objective is to translate solar X-ray telemetry into a measurable and deployable
software system for early detection, short-horizon forecasting, and impact assessment.

STELLA (Solar Temporal Event Learning \& Likelihood Assessment) combines temporal
solar observations, machine-learning-based event detection, forecasting, and
explainable risk assessment into one monitoring workflow. Rather than only identifying
a flare after it has developed, the system is designed to determine whether an event
is emerging, estimate its likelihood and severity, and communicate its potential
operational relevance.

% -----------------------------------------------------
% 2. TIME-TO-VALUE
% -----------------------------------------------------

\section{Time-to-value%
  \BvrHint{\ldots primary heuristic}}

We will deliver meaningful increments quickly by first establishing a working solar
telemetry pipeline and visualization layer, followed by a baseline flare detector.
Forecasting, explainability, and impact assessment will then be added incrementally.

We will prioritize fast feedback over premature model optimization by validating each
stage independently on historical solar observations before integrating the complete
pipeline. The separation of nowcasting and forecasting also allows the first iteration
to provide useful detection capability without waiting for the complete predictive
system.

Success is defined by what can be delivered and measured in the first iteration:
reliable telemetry processing, meaningful flare detection, measurable warning lead
time, and an interpretable user-facing workflow.

% -----------------------------------------------------
% 3. PROBLEM STATEMENT
% -----------------------------------------------------

\section{Problem Statement}

Solar flares are sudden releases of energy from the Sun that produce enhanced
electromagnetic radiation and can contribute to disruptions in radio communication,
navigation, satellite operations, and other space-dependent infrastructure. The
engineering challenge is not simply to identify a flare after it occurs, but to
recognize its emergence early enough to support preparation and response.

Several practical limitations make this difficult. Solar observations from different
instruments capture complementary physical characteristics, but are often analyzed
independently. Aditya-L1 provides valuable mission-specific observations, yet its
available history is much smaller than long-running solar observation archives.
Furthermore, solar background activity varies considerably over time, making fixed
thresholds prone to false alarms. Finally, a prediction is less useful operationally
when users cannot understand why it was generated or what systems may be affected.

\textbf{Why this matters now (contextual relevance):} Increasing dependence on
satellite communication, navigation, remote sensing, and space-based infrastructure
makes timely space-weather intelligence increasingly valuable. A system that turns
raw solar observations into understandable warnings can reduce the gap between
scientific monitoring and operational decision-making.

% -----------------------------------------------------
% 4. PROPOSED SOLUTION
% -----------------------------------------------------

\section{Proposed Solution%
  \BvrHint{\ldots Software Proposition}}

\subsection{Overview}

Build a real-time solar intelligence system that combines complementary X-ray
observations associated with Aditya-L1 with temporal machine learning and operational
visualization. STELLA will use SoLEXS soft X-ray observations and HEL1OS hard X-ray
observations to capture complementary thermal and higher-energy signatures of solar
activity.

The system will perform telemetry ingestion, preprocessing, adaptive anomaly
detection, temporal feature extraction, flare nowcasting, short-horizon forecasting,
explainable prediction, and impact visualization. The final interface will present
the current solar state, event likelihood, predicted flare severity, confidence,
warning lead time, and potential infrastructure risk.

\subsection{Core workflow}

The core workflow is intentionally simple:

\begin{enumerate}
    \item Solar X-ray telemetry is ingested and synchronized.
    \item Temporal and spectral features are extracted and unusual activity is
    identified using adaptive statistics.
    \item A Conv1D model performs short-window flare nowcasting.
    \item A dilated Temporal Convolutional Network evaluates longer temporal context
    and estimates future event likelihood and severity.
    \item The system generates an explanation and maps the prediction to potential
    infrastructure risks before presenting an alert through the dashboard.
\end{enumerate}

This creates a complete operational chain:

\begin{center}
\textbf{Telemetry $\rightarrow$ Detection $\rightarrow$ Forecast $\rightarrow$
Explanation $\rightarrow$ Impact $\rightarrow$ Alert}
\end{center}

\subsection{Operational constraints}

The system is designed as an engineering prototype rather than a replacement for
operational space-weather agencies. Historical solar observations will be used for
training and validation where mission-specific data is limited, while historical
NOAA GOES observations can provide a larger training archive for transfer learning.

Where continuous real-time telemetry is unavailable, historical replay will reproduce
the chronological behavior of an operational stream. Infrastructure outputs will be
presented as modeled risk indicators rather than deterministic predictions of physical
damage. The architecture will also remain modular so that additional solar instruments
and data products can be incorporated later.

% -----------------------------------------------------
% 5. SOLUTION APPROACH
% -----------------------------------------------------

\section{Solution Approach%
  \BvrHint{\ldots Engineering focus}}

This is an engineering course project, so the focus is on building a measurable,
deployable temporal intelligence pipeline rather than claiming novel fundamental
solar-physics research.

\subsection{Temporal feature and anomaly assistance%
  \BvrHint{\ldots physics-informed}}

STELLA will derive features including soft and hard X-ray flux, flux rise rate,
spectral hardness, rolling statistics, adaptive z-scores, and cross-instrument
temporal relationships. A key physics-informed feature is the relationship between
hard X-ray emission and subsequent soft X-ray evolution, motivated by the Neupert
effect. This relationship will be used as an additional predictive signal rather
than as a deterministic forecasting rule.

To handle changing solar background activity, the system will use rolling median
absolute deviation (MAD):

\[
MAD = median(|x_i - median(x)|)
\]

and an adaptive standardized score:

\[
Z_{adaptive} =
\frac{x-median(x)}
{1.4826\times MAD+\epsilon}
\]

This allows anomaly detection to adapt to changing background conditions rather
than relying on a single global threshold.

\subsection{Two-stage temporal architecture%
  \BvrHint{\ldots model design}}

The prediction pipeline separates immediate detection from longer-horizon forecasting.
The \textbf{Conv1D nowcaster} operates on a short temporal window and estimates
whether significant flare-like activity is emerging. The \textbf{dilated TCN
forecaster} then uses a longer temporal context to estimate event likelihood,
severity, and expected warning lead time.

This cascade is intended to keep the two objectives distinct: the first stage is
optimized for rapid recognition of emerging signatures, while the second captures
longer temporal dependencies required for forecasting.

\subsection{Transfer learning%
  \BvrHint{\ldots limited-data strategy}}

Because mission-specific training data is limited, historical NOAA GOES X-ray
observations will be used to learn generalized temporal representations. The model
will first be trained on the larger historical archive, then adapted using available
Aditya-L1 observations and evaluated on held-out events.

Chronological dataset splits will be preferred to random splits so that future
observations do not inadvertently influence training and evaluation. This also makes
the validation procedure more representative of the intended forecasting scenario.

\subsection{Web architecture%
  \BvrHint{\ldots deployable solution}}

The system will use a modular web architecture. The frontend will provide the
monitoring dashboard, while a FastAPI backend will handle telemetry ingestion,
inference requests, historical data, alerts, and system status. Python-based
inference services will execute the feature-processing and temporal models, with
REST APIs and WebSockets supporting standard requests and live updates.

The architecture separates frontend, backend, data processing, and inference so
that individual components can be tested and scaled independently.

\subsection{CI/CD and fast delivery enabler}

CI/CD will automatically build the application, run unit and integration tests,
validate API contracts and inference sanity checks, and produce repeatable staging
deployments. This supports frequent incremental releases while reducing the risk
of breaking the complete monitoring pipeline during model or interface changes.

% -----------------------------------------------------
% 6. EVALUATION CRITERION
% -----------------------------------------------------

\section{Evaluation Criterion%
  \BvrHint{\ldots measurable and attributable}}

We define success using measurable metrics tied directly to the forecasting workflow.
The primary metric is \textbf{warning lead time}: the time between the first valid
system warning and the defined reference time for the corresponding solar event.

\subsection{Primary evaluation metric}

\textbf{Warning lead time:} The initial target is at least 15 minutes of useful
warning for qualifying events, with approximately 30 minutes as a stretch objective.
The value will be calculated using timestamped telemetry, model predictions, and
historical flare-event catalogues.

\subsection{Secondary evaluation metrics}

\begin{itemize}
    \item \textbf{Probability of Detection (POD):} fraction of qualifying events
    correctly identified.
    \item \textbf{False Alarm Ratio (FAR):} fraction of generated warnings that do
    not correspond to qualifying events.
    \item \textbf{Critical Success Index (CSI):} combined measure of hits, misses,
    and false alarms.
    \item \textbf{Inference latency:} time required to convert new telemetry into
    a prediction.
    \item \textbf{Forecast calibration:} agreement between predicted likelihoods
    and observed event frequencies.
\end{itemize}

\subsection{Pilot validation plan%
  \BvrHint{\ldots high-level}}

Validation will proceed from offline evaluation to operational replay. First, models
will be evaluated on chronologically held-out historical events. Historical telemetry
will then be streamed through the complete pipeline to simulate real-time operation.
Finally, the deployed system will be evaluated for ingestion latency, inference
latency, warning generation, false alarms, and dashboard responsiveness.

% -----------------------------------------------------
% 7. SCALABILITY
% -----------------------------------------------------

\section{Scalability%
  \BvrHint{\ldots with foresight}}

The proposition should scale in both theoretical and practical senses.

\begin{enumerate}
    \item \textbf{Scalable (theoretically):} The system can support increasing
    telemetry volume and additional instruments through decoupled frontend, backend,
    and inference services, stateless APIs, asynchronous processing, and indexed
    historical data.
    
    \item \textbf{Operationally deployable (practically):} The system can run on
    common cloud or local infrastructure using containerized services, managed or
    local databases, standard Python ML environments, WebSocket-enabled APIs, and
    CI/CD pipelines.
\end{enumerate}

The modular design also allows new instruments or derived solar data products to be
introduced without redesigning the complete application.

% -----------------------------------------------------
% 8. ENGINE AVAILABILITY HEURISTIC
% -----------------------------------------------------

\section{Engine availability heuristic}

A mature set of ``engines'' is available to implement this as a solar intelligence
project. Historical solar X-ray archives provide training and validation data,
while Python-based scientific and machine-learning frameworks support temporal
feature processing and model development. Conv1D and Temporal Convolutional Network
architectures provide established mechanisms for temporal learning.

FastAPI can provide REST and WebSocket services, while React/Vite can support the
monitoring interface. Standard testing frameworks and CI runners can provide
repeatable validation and deployment.

We will use these mature components to focus on temporal feature fusion, forecasting,
explainability, system integration, and measurement rather than inventing new
infrastructure.

% -----------------------------------------------------
% 9. PROJECT SCOPE AND DELIVERABLES
% -----------------------------------------------------

\section{Project scope and deliverables%
  \BvrHint{\ldots iteration-friendly}}

\subsection{Initial deliverable%
  \BvrHint{\ldots first iteration}}

The first iteration will establish the minimum end-to-end system:

\begin{itemize}
    \item Solar telemetry ingestion and preprocessing.
    \item Historical GOES data pipeline and temporal feature extraction.
    \item Adaptive MAD-based anomaly detection.
    \item Initial Conv1D nowcasting model.
    \item Basic inference API and telemetry visualization.
    \item Prediction timestamp logging, CI testing, and staging deployment.
\end{itemize}

The objective is to demonstrate a complete working path from telemetry to an
interpretable detection output before adding the more advanced forecasting and
impact layers.

\subsection{Subsequent deliverables}

Subsequent iterations will add the dilated TCN forecasting model, SoLEXS/HEL1OS
multi-instrument fusion, GOES-to-Aditya-L1 transfer learning, lead-time estimation,
explainable AI, historical event catalogues, and the India-specific impact layer.

The final dashboard will combine live telemetry, event likelihood, predicted class,
confidence, lead time, explanation, alerts, and infrastructure risk into a single
operational view.

% -----------------------------------------------------
% 10. RISKS AND MITIGATIONS
% -----------------------------------------------------

\section{Risks and mitigations}

\begin{itemize}
    \item \textbf{Limited mission-specific data:} use historical GOES observations
    for pre-training and adapt using available Aditya-L1 data.
    
    \item \textbf{Instrument domain shift:} apply instrument-specific normalization
    and validate individual signals before fusion.
    
    \item \textbf{False alarms:} use adaptive thresholds and explicitly measure FAR
    under different solar activity levels.
    
    \item \textbf{Data leakage:} use chronological splits and prevent observations
    from the same event from crossing training and testing partitions.
    
    \item \textbf{Limited real-time telemetry:} maintain a chronological historical
    replay mode for end-to-end testing.
    
    \item \textbf{Model overfitting:} compare against simple baselines and evaluate
    on held-out events.
    
    \item \textbf{Impact uncertainty:} communicate infrastructure results as modeled
    risk indicators rather than deterministic damage predictions.
\end{itemize}

% -----------------------------------------------------
% 11. SUMMARY
% -----------------------------------------------------

\section{Summary}

This project proposes a deployable solar intelligence platform that converts temporal
X-ray observations into early-warning information for significant solar events. It
emphasizes time-to-value through rapid iterations, CI/CD for reliable releases, and
measurable evaluation through warning lead time, detection performance, false alarms,
and inference latency.

STELLA combines multi-instrument observations, adaptive anomaly detection, Conv1D
nowcasting, TCN forecasting, transfer learning, explainable AI, and impact
visualization within a single workflow. The system remains focused on engineering
a reliable decision-support pipeline rather than claiming novel solar-physics
discovery.

Its modular architecture supports future instruments and data products, while
historical replay enables meaningful validation even when continuous real-time
telemetry is unavailable. The final system provides a unified interface to observe
solar activity, identify emerging events, estimate their likelihood, understand
prediction drivers, and assess potential infrastructure relevance.

\end{document}

